\documentclass[10pt,a4paper]{amsart}
\usepackage{amsmath,amssymb,amsthm,graphicx}
\usepackage{verbatim}
\newenvironment{reference}{\comment}{\endcomment}
%\newenvironment{reference}{}{}
\newenvironment{slogan}{\comment}{\endcomment}
\newenvironment{history}{\comment}{\endcomment}
% For commutative diagrams we use Xy-pic
\usepackage[all]{xy}

% We use 2cell for 2-commutative diagrams.
\xyoption{2cell}
% Theorem environments.
%
\theoremstyle{plain}
\newtheorem{theorem}[subsection]{Theorem}
\newtheorem{proposition}[subsection]{Proposition}
\newtheorem{lemma}[subsection]{Lemma}

\theoremstyle{definition}
\newtheorem{definition}[subsection]{Definition}
\newtheorem{example}[subsection]{Example}
\newtheorem{exercise}[subsection]{Exercise}
\newtheorem{situation}[subsection]{Situation}

\theoremstyle{remark}
\newtheorem{remark}[subsection]{Remark}
\newtheorem{remarks}[subsection]{Remarks}

\numberwithin{equation}{subsection}

% Macros
%
\def\lim{\mathop{\mathrm{lim}}\nolimits}
\def\colim{\mathop{\mathrm{colim}}\nolimits}
\def\Spec{\mathop{\mathrm{Spec}}}
\def\Hom{\mathop{\mathrm{Hom}}\nolimits}
\def\Ext{\mathop{\mathrm{Ext}}\nolimits}
\def\SheafHom{\mathop{\mathcal{H}\!\mathit{om}}\nolimits}
\def\SheafExt{\mathop{\mathcal{E}\!\mathit{xt}}\nolimits}
\def\Sch{\mathit{Sch}}
\def\Mor{\mathop{\mathrm{Mor}}\nolimits}
\def\Ob{\mathop{\mathrm{Ob}}\nolimits}
\def\Sh{\mathop{\mathit{Sh}}\nolimits}
\def\NL{\mathop{N\!L}\nolimits}
\def\CH{\mathop{\mathrm{CH}}\nolimits}
\def\proetale{{pro\text{-}\acute{e}tale}}
\def\etale{{\acute{e}tale}}
\def\QCoh{\mathit{QCoh}}
\def\Ker{\mathop{\mathrm{Ker}}}
\def\Im{\mathop{\mathrm{Im}}}
\def\Coker{\mathop{\mathrm{Coker}}}
\def\Coim{\mathop{\mathrm{Coim}}}

% Boxtimes
%
\DeclareMathSymbol{\boxtimes}{\mathbin}{AMSa}{"02}

%
% Macros for moduli stacks/spaces
%
\def\QCohstack{\mathcal{QC}\!\mathit{oh}}
\def\Cohstack{\mathcal{C}\!\mathit{oh}}
\def\Spacesstack{\mathcal{S}\!\mathit{paces}}
\def\Quotfunctor{\mathrm{Quot}}
\def\Hilbfunctor{\mathrm{Hilb}}
\def\Curvesstack{\mathcal{C}\!\mathit{urves}}
\def\Polarizedstack{\mathcal{P}\!\mathit{olarized}}
\def\Complexesstack{\mathcal{C}\!\mathit{omplexes}}
% \Pic is the operator that assigns to X its picard group, usage \Pic(X)
% \Picardstack_{X/B} denotes the Picard stack of X over B
% \Picardfunctor_{X/B} denotes the Picard functor of X over B
\def\Pic{\mathop{\mathrm{Pic}}\nolimits}
\def\Picardstack{\mathcal{P}\!\mathit{ic}}
\def\Picardfunctor{\mathrm{Pic}}
\def\Deformationcategory{\mathcal{D}\!\mathit{ef}}
\usepackage{cleveref}
\crefname{originallemma}{original lemma}{original lemmas}
\Crefname{originallemma}{Original lemma}{Original lemmas}
\makeatletter
\newlabel{divisors-lemma-quasi-projective-Noetherian-pic-effective-Cartier}{{0AYM}{source}{}{sourcepointer-1}{}}
\newlabel{divisors-lemma-quasi-projective-Noetherian-pic-effective-Cartier@cref}{{[originallemma][0][]0AYM}{[1][0][]source}{}{}{}}
\makeatother
\crefname{originallemma}{original lemma}{original lemmas}
\Crefname{originallemma}{Original lemma}{Original lemmas}
\makeatletter
\newlabel{lemma-gysin-bivariant}{{0B79}{source}{}{sourcepointer-2}{}}
\newlabel{lemma-gysin-bivariant@cref}{{[originallemma][0][]0B79}{[1][0][]source}{}{}{}}
\makeatother
\crefname{originallemma}{original lemma}{original lemmas}
\Crefname{originallemma}{Original lemma}{Original lemmas}
\makeatletter
\newlabel{lemma-push-proper-bivariant}{{0EPK}{source}{}{sourcepointer-3}{}}
\newlabel{lemma-push-proper-bivariant@cref}{{[originallemma][0][]0EPK}{[1][0][]source}{}{}{}}
\makeatother
\crefname{originallemma}{original lemma}{original lemmas}
\Crefname{originallemma}{Original lemma}{Original lemmas}
\makeatletter
\newlabel{lemma-support-cap-effective-Cartier}{{02T9}{source}{}{sourcepointer-4}{}}
\newlabel{lemma-support-cap-effective-Cartier@cref}{{[originallemma][0][]02T9}{[1][0][]source}{}{}{}}
\makeatother
\crefname{originalreference}{original reference}{original references}
\Crefname{originalreference}{Original reference}{Original references}
\makeatletter
\newlabel{situation-setup}{{02QL}{source}{}{sourcepointer-5}{}}
\newlabel{situation-setup@cref}{{[originalreference][0][]02QL}{[1][0][]source}{}{}{}}
\makeatother
\title{Chern classes of d + 1 invertible sheaves on a finite type scheme with an ample invertible sheaf vanish in degree d + 1}
\author{The Stacks Project Authors}
\date{stacks-project:chow.tex:local-snapshot-20261012; GNU FDL 1.2 or later}
\begin{document}
\maketitle
\noindent\textit{Editorial note.} All delivered mathematical text is the author's own, copied from the pinned Stacks Project chapter file chow.tex; no statement or proof step was rewritten, reordered or removed.\par\medskip
\noindent\textit{Editorial note.} Preamble provenance: the pinned chapter file opens with \textbackslash{}input\{preamble\} and carries no preamble text of its own, so the project's single Stacks-wide preamble staged at sources/stacks-chow/preamble.tex is excerpted verbatim as three spans, lines 8-12 (the verbatim package together with the project's own reference, slogan and history display-environment definitions, including the author's commented-out older form of the reference environment), lines 14-18 (Xy-pic setup) and lines 154-217 (the native theorem declarations and mathematical macros). The lines 8-12 span is required because a counted statement or its surrounding text may use those display environments, which the project typesets as comments. Nothing was added, removed or redefined; the wrapper class and the amsmath/amssymb/amsthm/graphicx base packages are supplied by the assembler.\par\medskip
\noindent\textit{Editorial note.} Same-chapter pointers are located at their \textbackslash{}label lines inside the pinned chapter file; cross-chapter pointers are located in the project's own tag list staged in the snapshot (sources/stacks-chow/tags.txt), at the verbatim lines carrying their canonical tags.\par\medskip
\noindent\textit{Allowed standard prerequisites (editorial).}\begin{itemize}
\item Every invertible sheaf on a quasi-projective Noetherian scheme is a difference of effective Cartier divisors. Divisors, label divisors-lemma-quasi-projective-Noetherian-pic-effective-Cartier, invoked at proof line 6389. Writes the last invertible sheaf as the difference of two effective Cartier divisors, so that by additivity one may assume it is the sheaf of a single effective Cartier divisor.
\item The Gysin bivariant class of an effective Cartier divisor and the expression of its first Chern class. Same chapter chow.tex, labels lemma-gysin-bivariant (proof line 6400) and lemma-support-cap-effective-Cartier (proof line 6402). Provides the bivariant class attached to the inclusion of the divisor and identifies the composite of pushforward with it as the first Chern class of the divisor's sheaf.
\item Pushforward of bivariant classes and the definition of a bivariant class. Same chapter chow.tex, label lemma-push-proper-bivariant, invoked at proof line 6403. Makes sense of the composite of a pushforward with a bivariant class, which is what allows the product of the Chern classes to be rewritten as a pushforward of a product on the divisor.
\end{itemize}\medskip
\ifdefined\hypertarget\hypertarget{sourcepointer-1}{}\fi
\noindent\textit{Source reference.} Original lemma 0AYM: sources/stacks-chow/tags.txt, source lines 13795--13795. Navigation pointer for \textbackslash{}ref\{divisors-lemma-quasi-projective-Noetherian-pic-effective-Cartier\} used inside the counted span (line 6389); the argument that consumes it is retained in the span.\par\medskip
\ifdefined\hypertarget\hypertarget{sourcepointer-2}{}\fi
\noindent\textit{Source reference.} Original lemma 0B79: sources/stacks-chow/pinned-original-source.tex, source lines 6149--6149. Navigation pointer for \textbackslash{}ref\{lemma-gysin-bivariant\} used inside the counted span (line 6400); the argument that consumes it is retained in the span.\par\medskip
\ifdefined\hypertarget\hypertarget{sourcepointer-3}{}\fi
\noindent\textit{Source reference.} Original lemma 0EPK: sources/stacks-chow/pinned-original-source.tex, source lines 6167--6167. Navigation pointer for \textbackslash{}ref\{lemma-push-proper-bivariant\} used inside the counted span (line 6403); the argument that consumes it is retained in the span.\par\medskip
\ifdefined\hypertarget\hypertarget{sourcepointer-4}{}\fi
\noindent\textit{Source reference.} Original lemma 02T9: sources/stacks-chow/pinned-original-source.tex, source lines 5138--5138. Navigation pointer for \textbackslash{}ref\{lemma-support-cap-effective-Cartier\} used inside the counted span (line 6402); the argument that consumes it is retained in the span.\par\medskip
\ifdefined\hypertarget\hypertarget{sourcepointer-5}{}\fi
\noindent\textit{Source reference.} Original reference 02QL: sources/stacks-chow/pinned-original-source.tex, source lines 1372--1372. Navigation pointer for \textbackslash{}ref\{situation-setup\} used inside the counted span (line 6375); the argument that consumes it is retained in the span.\par\medskip
\begin{lemma}
\label{lemma-vanish-above-dimension}
Let $(S, \delta)$ be as in Situation \ref{situation-setup}. Let $X$ be a
finite type scheme over $S$ which has an ample invertible sheaf.
Assume $d = \dim(X) < \infty$ (here we really mean dimension and
not $\delta$-dimension).
Then for any invertible sheaves $\mathcal{L}_1, \ldots, \mathcal{L}_{d + 1}$
on $X$ we have
$c_1(\mathcal{L}_1) \circ \ldots \circ c_1(\mathcal{L}_{d + 1}) = 0$
in $A^{d + 1}(X)$.
\end{lemma}

\begin{proof}
We prove this by induction on $d$. The base case $d = 0$ is true because
in this case $X$ is a finite set of closed points and hence every invertible
module is trivial. Assume $d > 0$. By Divisors, Lemma
\ref{divisors-lemma-quasi-projective-Noetherian-pic-effective-Cartier}
we can write $\mathcal{L}_{d + 1} \cong \mathcal{O}_X(D) \otimes
\mathcal{O}_X(D')^{\otimes -1}$ for some effective Cartier divisors
$D, D' \subset X$. Then $c_1(\mathcal{L}_{d + 1})$ is the difference
of $c_1(\mathcal{O}_X(D))$ and $c_1(\mathcal{O}_X(D'))$ and hence
we may assume $\mathcal{L}_{d + 1} = \mathcal{O}_X(D)$ for some
effective Cartier divisor.

\medskip\noindent
Denote $i : D \to X$ the inclusion morphism and denote
$i^* \in A^1(D \to X)$ the bivariant class given by the
Gysin hommomorphism as in Lemma \ref{lemma-gysin-bivariant}.
We have $i_* \circ i^* = c_1(\mathcal{L}_{d + 1})$
in $A^1(X)$ by Lemma \ref{lemma-support-cap-effective-Cartier}
(and Lemma \ref{lemma-push-proper-bivariant}
to make sense of the left hand side).
Since $c_1(\mathcal{L}_i)$ commutes with
both $i_*$ and $i^*$ (by definition of bivariant classes)
we conclude that
$$
c_1(\mathcal{L}_1) \circ \ldots \circ c_1(\mathcal{L}_{d + 1}) =
i_* \circ c_1(\mathcal{L}_1) \circ \ldots \circ c_1(\mathcal{L}_d) \circ i^* =
i_* \circ c_1(\mathcal{L}_1|_D) \circ \ldots \circ c_1(\mathcal{L}_d|_D)
\circ i^*
$$
Thus we conclude by induction on $d$. Namely, we have $\dim(D) < d$
as none of the generic points of $X$ are in $D$.
\end{proof}

\end{document}
